% GENERATED FILE. Edit canonical lessons, then run npm run generate:books.
% canonical-source-hash: fnv1a64:0bd7ca78e65cb56d
% canonical-lessons: GU-C490-sasro, GU-C490-diyar, GU-C490-nanand, GU-C490-jethani, GU-C490-gharsansar

\chapter{My Husband's Family}
\label{ch:gu-my-husband-s-family}
\hlchaptermodality{490}

\begin{chapteropening}
\textbf{By the end of this chapter:} \emph{I can name my husband's relatives.}

Complete the last lesson of chapter 490: I can name my husband's relatives.
\end{chapteropening}

\section[\texorpdfstring{\gu{સસરો} (sasro)}{સસરો (sasro)}]{\gu{સસરો} (sasro) — a father-in-law}

\label{lesson:GU-C490-sasro}



\begin{quote}
Before the new one: say the Gujarati for to recall, then the Gujarati for sympathy.
\end{quote}

\subsection*{You'll want to know: \gu{સસરો}}
\textbf{\gu{સસરો}} — \emph{sasro} — \textquotedblleft{}a father-in-law\textquotedblright{}.

\begin{quote}
\gu{મારા} \gu{સસરા} \gu{રોજ} \gu{સવારે} \gu{મંદિરે} \gu{જાય} \gu{છે}.
\end{quote}

\begin{quote}
\emph{mārā sasrā roj savāre mandire jāy che.}
\end{quote}

\begin{quote}
My father-in-law goes to the temple every morning.
\end{quote}

A masculine noun.

\begin{cousinweb}
\textbf{\gu{સસરો}} (\emph{sasro}) is inherited from Sanskrit \emph{śvaśura}, \textquotedblleft{}husband's father\textquotedblright{}; CDIAL lists the Gujarati word (CDIAL 12753). One's own father-in-law is spoken of in the respectful plural, \textbf{\gu{મારા} \gu{સસરા}} (\emph{mārā sasrā}).
\end{cousinweb}

\begin{grammarlens}[title={What you've built}]
One more word for describing what you lived and felt.
\end{grammarlens}

\subsection*{Your turn}
\begin{itemize}
\raggedright
  \item \emph{Say it:} \emph{sasro}
  \item \emph{Say it:} \emph{sasro}, once more
  \item \emph{Say it:} \emph{mārā sasrā roj savāre mandire jāy che.} — My father-in-law goes to the temple every morning.
  \item \emph{Recall:} say the Gujarati for to recall, then the Gujarati for sympathy, then say \emph{sasro} again
\end{itemize}

\begin{tcolorbox}[breakable,colback=teal!4,colframe=teal!35!black,title={Before you move on}]
What does \gu{સસરો} mean? (\textquotedblleft{}A father-in-law\textquotedblright{}.) Say it once more.
\end{tcolorbox}
\section[\texorpdfstring{\gu{દિયર} (diyar)}{દિયર (diyar)}]{\gu{દિયર} (diyar) — a husband's younger brother}

\label{lesson:GU-C490-diyar}



\begin{quote}
Before the new one: say the Gujarati for sympathy, then the Gujarati for a father-in-law.
\end{quote}

\subsection*{You'll want to know: \gu{દિયર}}
\textbf{\gu{દિયર}} — \emph{diyar} — \textquotedblleft{}a husband's younger brother\textquotedblright{}.

\begin{quote}
\gu{મીરાનો} \gu{દિયર} \gu{શાળામાં} \gu{ભણે} \gu{છે}.
\end{quote}

\begin{quote}
\emph{Mīrāno diyar śāḷāmā̃ bhaṇe che.}
\end{quote}

\begin{quote}
Mira's husband's younger brother is at school.
\end{quote}

A masculine noun.

\begin{cousinweb}
\textbf{\gu{દિયર}} (\emph{diyar}) is inherited from Sanskrit \emph{dēvṛ}, \textquotedblleft{}husband's younger brother\textquotedblright{}, by way of \emph{dēvara-}; CDIAL lists the Gujarati word (CDIAL 6546).
\end{cousinweb}

\begin{grammarlens}[title={What you've built}]
One more word for describing what you lived and felt.
\end{grammarlens}

\subsection*{Your turn}
\begin{itemize}
\raggedright
  \item \emph{Say it:} \emph{diyar}
  \item \emph{Say it:} \emph{diyar}, once more
  \item \emph{Say it:} \emph{Mīrāno diyar śāḷāmā̃ bhaṇe che.} — Mira's husband's younger brother is at school.
  \item \emph{Recall:} say the Gujarati for sympathy, then the Gujarati for a father-in-law, then say \emph{diyar} again
\end{itemize}

\begin{tcolorbox}[breakable,colback=teal!4,colframe=teal!35!black,title={Before you move on}]
What does \gu{દિયર} mean? (\textquotedblleft{}A husband's younger brother\textquotedblright{}.) Say it once more.
\end{tcolorbox}
\section[\texorpdfstring{\gu{નણંદ} (naṇand)}{નણંદ (naṇand)}]{\gu{નણંદ} (naṇand) — a husband's sister}

\label{lesson:GU-C490-nanand}



\begin{quote}
Before the new one: say the Gujarati for a father-in-law, then the Gujarati for a husband's younger brother.
\end{quote}

\subsection*{You'll want to know: \gu{નણંદ}}
\textbf{\gu{નણંદ}} — \emph{naṇand} — \textquotedblleft{}a husband's sister\textquotedblright{}.

\begin{quote}
\gu{મારી} \gu{નણંદ} \gu{દિવાળીમાં} \gu{અમારા} \gu{ઘરે} \gu{આવશે}.
\end{quote}

\begin{quote}
\emph{mārī naṇand divāḷīmā̃ amārā ghare āvśe.}
\end{quote}

\begin{quote}
My husband's sister will come to our house at Diwali.
\end{quote}

A feminine noun.

\begin{cousinweb}
\textbf{\gu{નણંદ}} (\emph{naṇand}) is inherited from Sanskrit \emph{nanāndṛ}, \textquotedblleft{}husband's sister\textquotedblright{}; CDIAL lists the Gujarati \emph{naṇãd} (CDIAL 6946).
\end{cousinweb}

\begin{grammarlens}[title={What you've built}]
One more word for describing what you lived and felt.
\end{grammarlens}

\subsection*{Your turn}
\begin{itemize}
\raggedright
  \item \emph{Say it:} \emph{naṇand}
  \item \emph{Say it:} \emph{naṇand}, once more
  \item \emph{Say it:} \emph{mārī naṇand divāḷīmā̃ amārā ghare āvśe.} — My husband's sister will come to our house at Diwali.
  \item \emph{Recall:} say the Gujarati for a father-in-law, then the Gujarati for a husband's younger brother, then say \emph{naṇand} again
\end{itemize}

\begin{tcolorbox}[breakable,colback=teal!4,colframe=teal!35!black,title={Before you move on}]
What does \gu{નણંદ} mean? (\textquotedblleft{}A husband's sister\textquotedblright{}.) Say it once more.
\end{tcolorbox}
\section[\texorpdfstring{\gu{જેઠાણી} (jeṭhāṇī)}{જેઠાણી (jeṭhāṇī)}]{\gu{જેઠાણી} (jeṭhāṇī) — a husband's elder brother's wife}

\label{lesson:GU-C490-jethani}



\begin{quote}
Before the new one: say the Gujarati for a husband's younger brother, then the Gujarati for a husband's sister.
\end{quote}

\subsection*{You'll want to know: \gu{જેઠાણી}}
\textbf{\gu{જેઠાણી}} — \emph{jeṭhāṇī} — \textquotedblleft{}a husband's elder brother's wife\textquotedblright{}.

\begin{quote}
\gu{જેઠાણી} \gu{રોજ} \gu{મને} \gu{રસોડામાં} \gu{મદદ} \gu{કરે} \gu{છે}.
\end{quote}

\begin{quote}
\emph{jeṭhāṇī roj mane rasoḍāmā̃ madad kare che.}
\end{quote}

\begin{quote}
My husband's elder brother's wife helps me in the kitchen every day.
\end{quote}

A feminine noun.

\begin{cousinweb}
\textbf{\gu{જેઠાણી}} (\emph{jeṭhāṇī}) goes back to a reconstructed \emph{*jyēṣṭhajāni-}, \textquotedblleft{}senior wife, elder brother's wife\textquotedblright{}, built from \emph{jyēṣṭha}, \textquotedblleft{}eldest\textquotedblright{}, and \emph{jāni}, \textquotedblleft{}wife\textquotedblright{}; CDIAL lists the Gujarati word (CDIAL 5287).
\end{cousinweb}

\begin{grammarlens}[title={What you've built}]
One more word for describing what you lived and felt.
\end{grammarlens}

\subsection*{Your turn}
\begin{itemize}
\raggedright
  \item \emph{Say it:} \emph{jeṭhāṇī}
  \item \emph{Say it:} \emph{jeṭhāṇī}, once more
  \item \emph{Say it:} \emph{jeṭhāṇī roj mane rasoḍāmā̃ madad kare che.} — My husband's elder brother's wife helps me in the kitchen every day.
  \item \emph{Recall:} say the Gujarati for a husband's younger brother, then the Gujarati for a husband's sister, then say \emph{jeṭhāṇī} again
\end{itemize}

\begin{tcolorbox}[breakable,colback=teal!4,colframe=teal!35!black,title={Before you move on}]
What does \gu{જેઠાણી} mean? (\textquotedblleft{}A husband's elder brother's wife\textquotedblright{}.) Say it once more.
\end{tcolorbox}
\section[\texorpdfstring{\gu{ઘરસંસાર} (gharsansār)}{ઘરસંસાર (gharsansār)}]{\gu{ઘરસંસાર} (gharsansār) — married life; running a household}

\label{lesson:GU-C490-gharsansar}



\begin{quote}
Before the new one: say the Gujarati for a husband's sister, then the Gujarati for a husband's elder brother's wife.
\end{quote}

\subsection*{You'll want to know: \gu{ઘરસંસાર}}
\textbf{\gu{ઘરસંસાર}} — \emph{gharsansār} — \textquotedblleft{}married life; running a household\textquotedblright{}.

\begin{quote}
\gu{લગ્ન} \gu{પછી} \gu{મીરાનો} \gu{ઘરસંસાર} \gu{સારો} \gu{ચાલે} \gu{છે}.
\end{quote}

\begin{quote}
\emph{lagna pachī Mīrāno gharsansār sāro cāle che.}
\end{quote}

\begin{quote}
Mira's married life is going well.
\end{quote}

A masculine noun.

\begin{cousinweb}
\textbf{\gu{ઘરસંસાર}} (\emph{gharsansār}) is built from \textbf{\gu{ઘર}} (\emph{ghar}), \textquotedblleft{}home\textquotedblright{}, and \emph{sansār}, Sanskrit \emph{saṃsāra}, \textquotedblleft{}worldly existence\textquotedblright{} (Monier-Williams).
\end{cousinweb}

\begin{grammarlens}[title={What you've built}]
One more word for describing what you lived and felt.
\end{grammarlens}

\subsection*{Your turn}
\begin{itemize}
\raggedright
  \item \emph{Say it:} \emph{gharsansār}
  \item \emph{Say it:} \emph{gharsansār}, once more
  \item \emph{Say it:} \emph{lagna pachī Mīrāno gharsansār sāro cāle che.} — Mira's married life is going well.
  \item \emph{Recall:} say the Gujarati for a husband's sister, then the Gujarati for a husband's elder brother's wife, then say \emph{gharsansār} again
\end{itemize}

\begin{tcolorbox}[breakable,colback=teal!4,colframe=teal!35!black,title={Before you move on}]
What does \gu{ઘરસંસાર} mean? (\textquotedblleft{}Married life; running a household\textquotedblright{}.) Say it once more.
\end{tcolorbox}
